% magyar LaTeX article sablon 

\documentclass{article}
\usepackage[utf8]{inputenc}
\usepackage[T1]{fontenc}
\usepackage[magyar]{babel}
\usepackage{verbatim}
% margók beállítása
\usepackage{geometry} 
\geometry{left=1.5cm, right=1.5cm, top=1.5cm, bottom=1.5cm}

%opening
\title{README}
\author{Fábián Bernát}
\begin{document}
- a cím- és névadatokat tartalmazó szövegfájl:
\\
data/database.csv
\\\\
– a fölhasznált \LaTeX sablon:
\\
sample.tex
\\\\
– a körleveleket generáló szkriptek, programok:
\\
data/LetterGenerator.py, data/DataGenerator.py
\\\\
– a létrehozott, nyomtatásra kész (PDF) körlevelek:
\\
pdf\_outputs/*.pdf
\\\\
– a felhasznált \LaTeX csomagok, stílusfájlok (ha ezek közismertek,
könnyen elérhetők, elég felsorolni és az elérhetőségüket megadni.):
\begin{verbatim}
\documentclass[a4paper,12pt]{letter}
\usepackage[utf8]{inputenc}
\usepackage{lmodern}
\usepackage[magyar]{babel}
\usepackage{graphicx}
\usepackage{t1enc}
% margók beállítása
\usepackage{geometry} 
\geometry{left=2.5cm, right=2.5cm, top=2.5cm, bottom=2.5cm}
\usepackage{fancyhdr} % Fancy header/footer
\usepackage{enumerate}
\usepackage{hyperref}
\end{verbatim}
\end{document}